\section{Limitations and Next Work}
\label{sec:limitations}

The most important limitation is the sodium-load extrapolation. Every produced-water row is outside the original Rezaee DOE \NaIon/\LiIon training range, so the surrogate is currently appropriate for process screening and interface development rather than final design. A second limitation is the gap between the calibrated distribution surrogate and a fully accepted direct reactive ePC-SAFT LLE closure. The current package-facing evidence is valuable, but the manuscript should not claim an accepted physical split until the direct solve, reaction convention, and source-supported residuals are all closed.

The next paper-development steps are:
\begin{enumerate}
    \item Replace internal path references with stable citation, table, and artifact identifiers.
    \item Promote the final Rezaee regression table, residuals, and uncertainty diagnostics into the manuscript.
    \item Add the full \Prommis MSContactor equations and formal \IDAES costing assumptions.
    \item Add a transparent pretreatment model for divalent removal before the clean \LiIon/\NaIon extraction feed.
    \item Compile a final source-backed literature review covering produced-water lithium resources, hydrophobic \DES extraction, \TOPO co-extraction, and thermodynamic surrogate limits.
\end{enumerate}

\section{Conclusions}
\label{sec:conclusions}

This scaffold establishes the first manuscript structure for the produced-water lithium extraction case study. The active story is a Rezaee-calibrated \DES/\TOPO distribution surrogate applied to a pretreated \Smackover \LiIon/\NaIon stream, with explicit transfer variables for \Prommis/\IDAES process screening. The draft deliberately preserves the boundary between process-screening evidence and final predictive thermodynamics so that later revisions can strengthen the paper without overstating the current model.
